\section{Gramatika datalogu pre ANTLR}
Program ANTLR potrebuje na vygenerovanie parseru gramatiku datalogu. V tejto kapitole ju zostrojíme
a odôvodníme, prečo sme zvolili práve takúto gramatiku.

\vline

ANTLR používa iný zápis pravidiel gramatiky ako sme definovali v teoretickej časti. Napríklad
pravidlo $\sigma \rightarrow w_1 \, | \, w_2$ sa zapíše v ANTLR ako \verb$sigma: w1 | w2;$.
Terminály sa píšu do jednoduchých úvodzoviek, napríklad \verb$'('$. Tokeny sa píšu veľkými
písmenami, napríklad \verb$NUMBER$. Užitočným rozšírením je možnosť použiť zátvorky na zoskupovanie,
napríklad \verb$(w_1 | w_2)$, čo znamená, že sa môže použiť buď $w_1$ alebo $w_2$, \verb$(w)*$
znamená, že $w$ sa môže použiť nula alebo viac krát a \verb$(w)+$ znamená, že $w$ se musí použiť
aspoň raz. Viacej o zápise gramatiky v ANTLR sa dočítate v dokumentácii \cite{antlrDoc}.

\vline

Datalogový program (teória s dotazom) sa skladá z niekoľkých klauzúl (niektoré z nich môžy byť
faktami) nasledované dotazom. Do gramatiky (pre ANTLR) to prepíšeme ako 
\begin{verbatim}
program: (clause | definition)* query;
\end{verbatim}
kde \verb|program| bude počiatočný neterminál celej gramatiky, \verb|clause|, resp.
\verb|definition|, resp. \verb|query| bude neterminál, z ktorého sa odvodí klauzula, resp.
definícia, resp. dotaz.

Neterminál \verb|definition| sa bude prepisovať priamočiaro. Najprv treba zadať meno predikátu, ktorý
definujeme a potom napísať $n$-tice termov, pre ktorý je platný.

\begin{verbatim}
definition: name BODY_HEAD_SEPARATOR '{' '}' EOL |
            name BODY_HEAD_SEPARATOR '{'
                term_tuple (',' term_tuple)*
            '}' EOL;

term_tuple: '(' ')' |
            '(' term (',' term)* ')';
\end{verbatim}

Neterminál \verb|clause| bude trochu komplikovanejšia, pretože musí obsahovat hlavu a telo klauzuly. V tele
klauzuly sa môžu vyskytovať predikáty, negované predikáty, ale aj vstavané predikáty a negované
vstavané predikáty.

\begin{verbatim}
clause: name '(' ')' BODY_HEAD_SEPARATOR 
                    (predicate | built_in_predicate) 
                    (',' (predicate | built_in_predicate))* EOL |
        name '(' term (',' term)* ')' BODY_HEAD_SEPARATOR 
                    (predicate | built_in_predicate)
                    (',' (predicate | built_in_predicate))* EOL;

predicate:          normal_predicate |
                    negative_predicate;
negative_predicate: NOT normal_predicate;

built_in_predicate:          normal_built_in_predicate |
                             negative_built_in_predicate;
negative_built_in_predicate: NOT normal_built_in_predicate;
\end{verbatim}

Normálny vstavaný predikát sa skladá z mena a z $n$-tice termov, rovnako ako hlava klauzuly. 

\begin{verbatim}
normal_predicate: name '(' ')' |
                  name '(' term (',' term)* ')';
\end{verbatim}

Vstavaný predikát je podobný normálnemu predikátu, až na to, že ho musíme označiť, že jeho parametre
nie sú iba termy, ale aj predikáty a vstavané predikáty, ktoré môžu byť rôzne usporiadané v zoznamoch.

\begin{verbatim}
normal_built_in_predicate: BUILT_IN name '(' ')'|
                           BUILT_IN name '(' 
                               parameter (',' parameter)* 
                           ')';

parameter: '[' ']' |
           '[' parameter (',' parameter)* ']' |
           normal_predicate |
           normal_built_in_predicate |
           term;
\end{verbatim}

Token \verb|BUILT_IN| by sme mohli vynechať. V takom prípade by ale kompilátor nevedel rozlíšiť
vstavané predikáty od normálnych predikátov. To by pri následnom preklade do relačnej algebry
znamenalo, že každý predikát, čo nebol definovaný faktom alebo klauzulou, by sa musel považovat za
vstavaný predikát. Keďže ale chceme aby kompilátor vedel fungovať samostatne, bez prístupu do
\textit{databázy}, nevieme aké vstavané predikáty existujú. Kvôli tomu by sme nevedeli vyhodiť
kompilačnú chybu, keď v klauzule použijeme nedefinovaný predikát.

\vline

Keďže povolujeme iba kladné dotazy, tak \verb|query| sa môže prepísať iba na normálny predikát,
obyčajný alebo vstavaný.

\begin{verbatim}
query: QUERY_MARKER (normal_predicate | normal_built_in_predicate) EOL;
\end{verbatim}

Použité tokeny sú definované nasledovne, pričom \verb|[0-9]| resp. \verb|[a-z]| znamená ľubovoľný
znak z množiny číslic resp. písmen:

\begin{minipage}{0.9\textwidth}
\begin{verbatim}
QUERY_MARKER: '?-';
BODY_HEAD_SEPARATOR: ':-';
SLC: '%';
MLC_START: '/*';
MLC_END: '*/';
EOL: '.';
NOT: '\\+';
BUILT_IN: '$';

UPPER: [A-Z]+;
LOWER: [a-z]+;
NUMBER: [0-9]+;
\end{verbatim}
\end{minipage}
\medskip

Tieto tokeny nemuseli byť definované, ale na miestach kde sú použité, by sa mohlo použiť ich
definíciu. Náš prístup ale dovoľuje používateľovi jednoducho zmeniť tieto tokeny. Napríklad, ak by
chcel pre $\leftarrow$ používať alternatívny zápis \verb|<-|, stačí zmeniť definíciu tokenu
\verb|BODY_HEAD_SEPARATOR = '<-';|.

\vline

Názvy predikátov a funkčných symbolov sa skladajú z malých písmen, čísiel a znaku \verb|_|. Premenné
budú rovnaké, ale budú sa skladať z veľkých písmen. Všetky názvy sa musia začať písmenom.

\begin{verbatim}
name: LOWER (LOWER | NUMBER |'_')*;
function: LOWER (LOWER | NUMBER | '_')*;
variable: UPPER (UPPER | NUMBER | '_')*;
\end{verbatim}

Na koniec ešte ANTLR nastavíme, aby preskakoval \textit{whitespace} znaky a podporoval
jednoriadkové aj viacriadkové komentáre.

\begin{verbatim}
WS: [ \t\r\n]+ -> skip;
COMMENT: SLC ~[\n\r]* ([\n\r] | EOF) -> channel(HIDDEN);
MULTILINE_COMMENT: MLC_START (MULTILINE_COMMENT | .)*? 
                   (MLC_END | EOF) -> channel(HIDDEN);
\end{verbatim}